\section{Design decisions and rejected alternatives}\label{sec:tradeoffs}

The preceding sections describe the machinery; this one defends it. Every load-bearing
choice in the part-tree design is collected here as a decision record: the decision, the
rationale that wins, and the alternatives that lose --- with the specific failure that
disqualifies each. The entries are grouped by theme (ownership, caching, mutation,
the motor, the error and event surface, concurrency), but each stands alone and is
grounded in the code it defends.

A pattern recurs throughout and is worth naming up front: wherever a correctness property
can be made \emph{unrepresentable} --- by ownership type, by access control, by a typed
verb --- the design prefers that to a runtime check, and prefers a runtime check to a
convention. The alternatives below mostly lose because they defend the same property at a
weaker rung of that ladder.

\subsection{Ownership and structure}\label{sec:tradeoffs-ownership}

\paragraph{Linear ownership.}
Every \code{PartNode} owns its \code{Part} and its children through
\cppinline|std::unique_ptr| (\srcf{core/model/PartsModel.h}); \code{PartFactory} produces
\cppinline|std::unique_ptr<Part>| (\src{core/model/parts/PartFactory.h}{60}); and
\code{PartsModel::attach} consumes one by value, so ownership transfer is visible in the
signature and enforced by the move. No \cppinline|std::shared_ptr<Part>| exists anywhere
in the production tree. The consequence is that a part has exactly one owner at all
times: double-attach, membership in two models, and post-attach aliasing are compile
errors, not audited conventions.

\begin{rejected}[Rejected --- shared-pointer interior]
Holding parts by \cppinline|std::shared_ptr| and handing copies to consumers loses on two
counts, both silent. First, a retained alias is an unrouted write path: the composite
mass-delta gate (\cref{sec:caching}) backstops a mass edit --- the live sum moves on the
next read --- but the gate keys on composite mass \emph{alone}
(\src{core/model/PartsModel.cpp}{122}), so an equal-mass tensor edit through an alias
(a same-mass, different-diameter motor is the canonical case) flies a stale composite
inertia tensor with no error, ever. Widening the gate to notice tensor edits would mean
recomputing the tensor in order to validate the tensor --- the cache doing the work it
exists to skip. The staleness class therefore cannot be closed by checking; it must be
closed structurally, by making the alias uncompilable. Second, one shared part becomes
attachable under two models: each composite counts the same mass, and an edit routed
through one model leaves the other's caches stale with no event. Linear ownership deletes
both failure classes at the type level.
\end{rejected}

\paragraph{A recursive owned node tree, not a flat arena.}
The tree is stored as it is shaped: each node holds a vector of owned children and a
parent back-pointer, and \code{Part::Id} --- process-unique, never reused, fresh on every
clone --- doubles as the O(1) index key (\srcf{core/model/PartsModel.h}).

\begin{rejected}[Rejected --- flat arena with generational ids]
Contiguous node records indexed by integer slot, parent indexes, generation counters, and
a cached preorder were evaluated as an alternative storage scheme. The arena has the
cleanest invalidation story on paper --- structural edits bump a generation exhaustively,
by construction --- but it loses three ways. (1)~Slot reuse makes a bare index silently
wrong after freelist recycling; every retained handle must carry a generation stamp to
dodge the ABA hazard, and the handles that matter most --- GUI selection and undo ---
are held across arbitrary edits. \code{Part::Id} as a never-reused counter is immune
outright: a stale id fails the O(1) lookup safely instead of resolving to a stranger.
(2)~A payload-free generation bump can never drive Qt's granular row operations, which
need $\{parent, row\}$ \emph{before} the edit (\cref{sec:gui}); the event stream the GUI
consumes carries exactly that payload (\src{core/model/PartsModel.h}{143}).
(3)~Contiguity pays only when cache misses dominate traversal, i.e.\ at node counts
orders of magnitude beyond a rocket's. A design is tens of parts; the hot-path walk
touches each node once per integrator stage regardless of layout. The arena's genuinely
good ideas --- typed mutation verbs, exhaustive invalidation on structural edits,
stamp-after-success --- appear in the tree design without the storage cost.
\end{rejected}

\subsection{Cache placement and the gate}\label{sec:tradeoffs-caching}

\paragraph{Per-node caches.}
Every \code{PartNode} carries its own composite-mass-properties cache and its own
resolved-placement cache (\srcf{core/model/PartsModel.h}), not just the root.

\begin{rejected}[Rejected --- root-only caching]
Caching only at the root halves the field count and nothing else. Every composite query
is defined on \code{PartNode} precisely so that \emph{any} node is a valid sub-assembly
root --- the API surface that stage composites (multi-stage designs are planned) and
detached-subtree preview consume. Root-only caching turns each such query into a full
recompute with nowhere to memoize. Meanwhile the invalidation logic is identical either
way: the dirty walk already marks every ancestor to the root
(\srcf{core/model/PartsModel.h}), so per-node caching adds no invalidation complexity ---
only a few hundred bytes of \code{mutable} state across the whole tree.
\end{rejected}

\paragraph{The mass-delta gate compares exactly.}
The composite CM/tensor cache rebuilds iff the node is structurally dirty or the live
mass sum differs from the stamp by exact \code{==}, with a NaN sentinel forcing the first
build (\src{core/model/PartsModel.cpp}{122}). The key is the composite mass itself ---
already computed every stage as the ODE divisor (\src{core/sim/Propagator.cpp}{41}), so
the gate costs nothing beyond a comparison. The division of labor is deliberate: the gate
covers \emph{time} variation (a burning motor), routing covers \emph{edits} --- because
an equal-mass edit is invisible to any mass-keyed gate, as
\cref{sec:tradeoffs-routing} develops.

\begin{rejected}[Rejected --- an epsilon tolerance on the gate]
An approximate comparison trades the bit-exact post-burnout freeze for a tunable with no
principled value. After burnout the live sum repeats bit-for-bit, so exact \code{==}
freezes the cache with zero false rebuilds; during the slow tail of a burn, per-stage
mass deltas shrink below any fixed epsilon, so an epsilon gate silently freezes the CM
and tensor \emph{while mass is still draining} --- a stale-cache class with no error and
no way to pick the threshold that avoids it.
\end{rejected}

\begin{rejected}[Rejected --- time-keyed caching]
Caching keyed on the query time $t$ assumes $t$ is monotone and unique. The adaptive
RK45 integrator violates both: stages evaluate at repeated intermediate times, and a
rejected step re-queries times already seen. A $t$-keyed cache therefore either thrashes
(rebuild per stage --- exactly the hot-path cost the cache exists to avoid) or serves a
value built \emph{before an edit} at the same $t$. Keying on the live mass sum is immune
to both, and doubles as a physical statement: for a fixed structure, the composite CM and
tensor change only when the mass distribution does.
\end{rejected}

\subsection{Routed mutation}\label{sec:tradeoffs-routing}

\paragraph{Verbs on the model, private setters on the leaf.}
Every post-attach write path is a \code{PartsModel} verb --- \code{setPartMass},
\code{setPartInertia}, \code{setLink}, \code{mutatePart}, and the motor verbs --- each of
which writes and then dirties the correct cache chain before returning. \code{Part}'s
setters are \code{private} with \code{PartsModel} as the sole friend
(\src{core/model/parts/Part.h}{133}), so no other code can reach them.
\code{MutateHint} declares whether an arbitrary edit is mass-only or geometric, so a mass
edit never re-runs the placement resolve and envelope sweep. \Cref{lst:tradeoffs-setlink}
shows the pattern at its most instructive point: a seat edit moves resolved poses, and
with them the composite CM and tensor, at possibly unchanged mass --- so the verb dirties
\emph{both} chains, because no gate downstream could detect the change.

\begin{listing}[htbp]
\begin{minted}{cpp}
bool PartsModel::setLink(part::Part::Id child, const part::StationLink& link)
{
    PartNode* n = find(child);
    if(n == nullptr || n == root_.get())
    {
        return false; // a root has no incoming edge
    }
    const Event e{Event::LinkChanged, child, n->parent_->id(), n->rowInParent()};
    emitEvent(e, true);
    n->link_ = link;
    // A seat edit is geometry: resolved poses move and the composite CM/tensor move with them --
    // the mass gate alone cannot see an equal-mass geometry change, so both chains dirty.
    n->parent_->markMassDirty();
    n->parent_->markPlacementDirty();
    emitEvent(e, false);
    return true;
}
\end{minted}
\caption{A routed verb writes, then dirties the chains the write actually
invalidates (\srccap{core/model/PartsModel.cpp}{431}).}
\label{lst:tradeoffs-setlink}
\end{listing}

\begin{keyinsight}[The gate is a backstop, not a router]
The mass-delta gate catches mass drift from any source, but it is structurally blind to
equal-mass edits --- a tensor write, a seat move, a same-mass motor swap. Routing is
therefore load-bearing correctness, not style: it is the only invalidation-complete
option, and linear ownership (\cref{sec:tradeoffs-ownership}) is what makes the routed
surface exhaustive by construction.
\end{keyinsight}

\begin{rejected}[Rejected --- an observer or back-pointer on \code{Part}]
Giving the leaf a pointer to its node so its setters self-invalidate reintroduces tree
knowledge into the class whose contract is precisely that it has none
(\cref{sec:leaf}): every detached part would need null-check discipline on every setter.
Worse, a public self-invalidating setter bypasses the model's event stream --- caches
would refresh but no \code{Mutated} event would fire, so the GUI's mass column goes stale
over a perfectly fresh cache. Invalidation and notification must travel together, and
only a model-level verb sees both.
\end{rejected}

\begin{rejected}[Rejected --- immutable parts, replace-on-edit]
Modeling every edit as clone-with-new-value plus node swap makes staleness
unrepresentable, but at the cost of identity: \code{Part::Id} is fresh per instance, so
every mass edit churns the id. GUI selection (keyed on
\code{QModelIndex::internalId}, \src{gui/RocketTreeView.cpp}{140}), the model's typed
motor borrow, and any held undo handle would all invalidate on a slider drag, and a
scalar edit would masquerade as a structural detach/attach with event noise to match.
Stable identity across in-place edits is exactly what mutable-behind-verbs buys.
\end{rejected}

\begin{rejected}[Rejected --- polled version counters]
A per-part version integer bumped by public setters, aggregated and compared on read,
turns every composite query into an $O(N)$ version sweep on the per-stage hot path ---
paid on every read, forever, to detect edits that arrive through a handful of call
sites. And a version bump says \emph{that} something changed, not \emph{what}: placement
work would re-run on mass edits, or a second per-aspect counter would reinvent
\code{MutateHint} with strictly more state.
\end{rejected}

\begin{rejected}[Rejected --- \code{protected} leaf setters]
Protected access lets any subclass re-export a public mutator that skips invalidation
and events; a motor type wrapping its own state is the obvious candidate. Private plus a
single named friend (\src{core/model/parts/Part.h}{133}) makes the compiler enforce the
seam against subclasses too.
\end{rejected}

\subsection{The motor}\label{sec:tradeoffs-motor}

\paragraph{Live-derived state.}
\code{part::Motor} owns its \code{MotorModel} by value and derives everything at query
time: \code{getMass(t)} and \code{getI()} read the wrapped model live
(\src{core/model/parts/Motor.h}{46}). No seeded copy of the motor's mass or tensor
exists anywhere, so \code{PartsModel::setMotor}'s swap path reduces to \emph{replace the
wrapped model, dirty the mass chain} (\srcf{core/model/PartsModel.cpp}) --- and the
stale-seed divergence class has no representation to occupy.

\begin{rejected}[Rejected --- seeded base-field copies]
Seeding the leaf's mass/tensor fields from the motor at construction and re-seeding on
swap creates a write path that must \emph{remember} to invalidate --- and the one swap
most likely to be forgotten is the one the mass gate can never catch: equal mass,
different diameter, stale tensor, no error. Live derivation deletes the state that could
go stale instead of testing for its staleness.
\end{rejected}

\begin{keyinsight}[Burn-dependent grain inertia is enabled, not gated out]
The live-derivation contract deliberately declines to pin the motor's tensor constant.
Propellant-grain inertia that tracks the burn is desirable future modeling, and when it
arrives \code{getI()} grows a time parameter mirroring \code{getMass(t)} --- the
composite walk already has $t$ in hand at the call site. A design that froze the tensor
by contract would buy nothing today and cost that feature later.
\end{keyinsight}

\paragraph{Persisted as a database reference.}
A design file records its motor as a common-name reference plus the seat link on its edge
(\src{core/model/DesignSerializer.cpp}{205}); load re-resolves the name against the motor
database, and a miss degrades to design-without-motor with a logged warning
(\src{core/model/DesignSerializer.cpp}{257}) rather than a failed load. A \code{Motor}
appearing as an ordinary tree \code{<part>} element is rejected outright
(\src{core/model/DesignSerializer.cpp}{155}), so the reference is the one channel.

\begin{rejected}[Rejected --- inlining the thrust curve]
Embedding the full curve in the design file creates two sources of truth that skew as
the database updates, bloats every file with data the database already owns, and leaves
the motor's identity ambiguous (no key to re-resolve corrections against). Motors are
per-flight database entities; the airframe file names one, it does not contain one ---
and a missing name must not brick the airframe, hence degrade-with-warning rather than
throw.
\end{rejected}

\paragraph{Zero-span nodes host nothing.}
The envelope sweep excludes any candidate host whose axial span is zero
(\src{core/model/parts/Placement.cpp}{183}); the \code{Motor} node, with the base
zero-length geometry (\srcf{core/model/parts/Part.h}), is the standing example. The
solid-host rule samples radii across the host's interior span; a point mass has no
interior stations to sample, and admitting it as a host would let it shadow the real
enclosing part --- an offender ``fitting'' against a point.

\begin{rejected}[Rejected --- a synthetic motor envelope]
Inventing a nonzero casing envelope for the \code{Motor} node so it can participate in
the sweep would collide with the mount geometry the airframe already models --- the same
volume claimed twice --- and would oblige every future point-mass-like part (ballast,
electronics) to carry fake geometry just to stay out of the sweep's way. Exclusion by
span is one comparison and states the physics plainly.
\end{rejected}

\subsection{Errors, events, and the empty design}\label{sec:tradeoffs-errors}

\paragraph{Typed \code{std::expected} errors.}
Structural verbs return their cause: \code{attach} yields
\code{std::expected<Id, AttachError>} with \code{NullPart} / \code{NoSuchParent} /
\code{DuplicateMotor}, and \code{detach} yields the detached subtree or
\code{NoSuchId} / \code{IsRoot} (\src{core/model/PartsModel.h}{136}). The CLI switches
on the enumerator and prints a precise message per cause (\src{cli/Repl.cpp}{955});
tests pin each enumerator directly. On success, \code{detach}'s value \emph{is} the
subtree --- the natural undo/clipboard token --- so the error channel and the payload
share one signature.

\begin{rejected}[Rejected --- \code{bool} returns or logged no-ops]
A verb that fails silently (or with only a log line) forces callers to \emph{probe} for
failure --- counting children before and after the call --- and to guess at the cause in
their own error text, which then diverges per call site. A boolean is barely better: it
answers ``did it work'' but not ``why not,'' so the CLI cannot distinguish a bad parent
id from a second motor, and a test cannot pin the rejection it intends to provoke.
\end{rejected}

\paragraph{Dual aboutTo/did events.}
Every mutation fires its \code{Event} twice --- once with \code{before == true}, once
after (\src{core/model/PartsModel.h}{143}) --- carrying
$\{kind, id, parentId, row\}$. This is the exact shape Qt's model contract needs:
\code{beginInsertRows}/\code{beginRemoveRows} must run \emph{before} the structure
changes, while the pre-edit rows are still queryable, and the matching \code{end*} call
runs after (\src{gui/RocketTreeView.cpp}{32}). The core stays Qt-free; it simply emits
an event shape a Qt adapter can consume losslessly (\cref{sec:gui}).

\begin{rejected}[Rejected --- a single post-hoc notification]
A lone ``something changed'' callback after the fact can never satisfy the
\code{begin*}/\code{end*} protocol, so a Qt view built on it must fall back to
\code{beginResetModel}/\code{endResetModel} on every edit --- destroying selection and
expansion state on each keystroke of a mass field. The design keeps a coarse post-hoc
callback \emph{as well} (\srcf{core/model/RocketModel.h}) for consumers that only
repaint; the dual stream is a superset, not a compromise.
\end{rejected}

\paragraph{A nullable root.}
``No design'' is a real state and is represented directly: \code{hasDesign()} is a null
test on the root (\src{core/model/PartsModel.h}{161}). Every consumer is forced to
confront the empty state explicitly, which is what makes the CLI's contract crisp:
\code{ERR addpart: no design (use newdesign first)} (\src{cli/Repl.cpp}{906}), and
likewise for \code{checkdesign} and \code{listparts}.

\begin{rejected}[Rejected --- a placeholder root part]
A mandatory synthetic root makes the empty state unrepresentable rather than handled:
\code{hasDesign()} degrades to recognizing a magic part, the serializer must persist or
special-case it, the composite walk and the sweep process a part that is not part of the
rocket, and the CLI's verified \code{no design} error contract dissolves into
heuristics. A null pointer is one comparison and cannot be mistaken for geometry.
\end{rejected}

\subsection{Concurrency: single-writer now, snapshots later}\label{sec:tradeoffs-threading}

The model is single-writer by contract, and the contract is stated where it binds: in the
\code{PartsModel} class documentation and at the \code{PartNode} cache fields
(\srcf{core/model/PartsModel.h}). The \code{const} composite accessors are
\emph{reads that write} --- they rebuild and re-stamp \code{mutable} caches --- and even
\code{thrust(t)} writes a \code{mutable} log-once burnout latch inside a \code{const}
call chain (\src{core/model/MotorModel.h}{342}). Concurrent access, all-\code{const}
included, is therefore a data race by design, and the signatures must not be read as a
thread-safety promise.

\begin{rejected}[Rejected --- a lock or mutex]
The reads-that-write pattern defeats reader-writer locking outright: any read may need
to become the writer, so a shared/exclusive lock degenerates to exclusive on every
composite read --- a mutex acquisition per RK45 stage, on the hottest path in the
program (the propagator divides by \code{getMass(t)} every stage,
\src{core/sim/Propagator.cpp}{41}). Locking also fails the consumers it claims to serve:
a CM and a tensor are physically meaningful only as a pair from the same rebuild, and a
per-call mutex gives no such multi-value consistency unless held across grouped reads
--- inviting GUI-holds-lock stutter and callback re-entrancy deadlocks. Its one genuine
advantage, zero-staleness reads, is worthless to a GUI repainting from millisecond-old
frames.
\end{rejected}

\begin{designrule}[If threading arrives: snapshot, do not share]
The direction of record for a future worker-thread simulation is snapshotting: fly a
\code{PartNode::clone()} of the design on the worker --- edits during a flight then
apply to the \emph{next} launch, the sane semantics regardless of threading --- and
publish telemetry as immutable \code{StateData} frames. Batch work (motor-ladder sweeps,
Monte Carlo) takes one clone per flight and is embarrassingly parallel. Each thread owns
its data; only the handoff synchronizes. The one known design obligation for that day:
\code{Part::Id} is fresh on clone, so mapping results back onto the edited tree needs a
clone-time id correspondence. Until a real concurrency requirement exists, the
single-writer contract stays documented at the cache sites so threading work reopens it
deliberately instead of trusting \code{const}.
\end{designrule}

Taken together, the record reads as one decision made twelve ways: put each invariant
inside the smallest boundary that can actually defend it --- the type system where
possible, a friend seam or a typed verb where not --- and spend runtime checks only on
what structure cannot express. \Cref{sec:verification} closes the loop by showing how
the test suite pins each of these choices in place.
